\chapter{Methodik}
\label{chap:methodik}

In diesem Kapitel wird der Prozess der Datenakquise und -aufbereitung detailliert beschrieben. Da für die Fragestellung keine modernen digitalen Datensätze zur Verfügung standen, musste eine Brücke zwischen historischer Dokumentation und moderner Datenanalyse geschlagen werden.

\section*{Beschreibung der Datenquelle}
Als primäre Datenquelle dient das Werk \textit{"Life and Labour of the People in London"} von Charles Booth (1886--1903). Die Analyse stützt sich auf die detaillierten Tabellen, in denen die Bevölkerung Londons nach Berufsfeldern (section*s) und sozialen Klassen (A--H) kategorisiert wurde. Jede Zeile der Originaltabellen enthält Informationen über:
\begin{itemize}
    \item Die \textbf{Berufsklasse} (z. B. Labour, Artisans, Dealers).
    \item Demografische Merkmale wie \textbf{Heads of Families}, \textbf{Wives} und \textbf{Children} unter 15 Jahren.
    \item Die Verteilung der Personen auf die \textbf{sozialen Klassen A bis H}.
\end{itemize}

\section*{Problematik unklarer oder ungenauer Daten}
Der Datensatz beruht auf den gesammelten Daten von 'School Board Visitors' die Hausbesuche bei Familien mit Kindern im Schulalter und 2-3 Jahre vor dem Schulalter machten und den Beruf des Familienhaupts festhielten. Die Menge der in diesem Datensatz erfassten Menschen beläuft sich auf die Hälfte der Bevölkerung der beschriebenen Distrikte und ist davon ausgehend hochgerechnet.
Booth macht zunächst einige grundlegende Annahmen:
Er geht davon aus, dass Schulkinder in der Regel ältere und jüngere Geschwister haben und schätzt auf die gezählte Anzahl der Schulkinder, noch einmal die selbe Menge von Geschwisterkindern ("Children and young persons 13-20") zu den selben Klassen und Berufsgruppen hinzu.
Auch geht er davon aus, dass die verheirateten Männer verschiedener Klassen und Berufsgruppen mit Schulkindern generell repräsentativ sind für verheiratete Männer derselben Klasse und Beruf ohne Schulkinder, mit der Begründungen, sie würden ihren Beruf wählen bevor sie Kinder im Schulalter haben und auch nachdem die Kinder nicht mehr in der Schule sind weiterhin in diesem Beruf tätig sind.

Diese Schätzungen beinhalten keinerlei Überlegung, dass gerade ärmere Familien sich möglicherweise gegen (weitere) Kinder entschieden, da sie sich keine (weiteren) Kinder leisten konnten.
Abtreibungen waren in England seid 1803 illegal (Malicious Shooting or Stabbing Act 1803), doch Abtreibungsmittel wurden trotzdem großflächig in Frauenmagazinen beworben, mit vagen oder blumigen Umschreibungen als "Catholic Pills", "Womans Helper" oder auch als Abführmittel. Der Rückgang der Geburtenrate in Großbritannien von 35.5/1000 in den 1870ern zu 29/1000 in 1900, unterstützt die Interpretation, dass Frauen in der UK in diesem Zeitraum, die ihnen zur Verfügung stehenden Mittel nutzten um Schwangerschaften zu planen.

Es lässt sich anmerken, dass Booths gelieferten Daten teils ungenau sind. So werden an einigen Stellen Schätzungen vorgenommen, die nicht als solche ausgezeichnet sind. Da sich in den Datensätzen Werte befinden die offensichtlich Abschätzungen sind, wie z.B. eine Menge an exakt 40.000 Menschen in einer Kategorie.

Auch sind manche Formulierungen für uns heute nicht mehr nachvollziehbar. Was genau die Benennung einer Gruppe damals beinhaltet hat, ob es eine allgemein bekannte Definition gab usw.

Trotz all diesen Einwänden und Schätzungen, die Menge der erhobenen Daten ist mit 50\% der Bevölkerung der erfassten Stadtteile ein wertvoller Einblick in die damaligen Bevölkerungsstrukturen und vertretenen Berufsgruppen in London.

\section*{Datenerhebung und manuelle Transkription}
Die Datenerhebung erfolgte durch die manuelle Übertragung der Informationen aus digitalen Scans der Originalseiten in ein maschinenlesbares Format (CSV). Dieser Prozess umfasste:
\begin{itemize}
    \item \textbf{Digitalisierung:} Übertragung der gedruckten Tabellen in CSV-Dateien für jeden untersuchten Borough.
    \item \textbf{Qualitätssicherung:} Abgleich der übertragenen Werte mit den Originalwerten zur Vermeidung von Übertragungsfehlern.
    \item \textbf{Umgang mit historischen Begriffen:} Beibehaltung der originalen Bezeichnungen, um die historische Authentizität zu wahren.
\end{itemize}


\section*{Eingesetzte Tools und Infrastruktur}
Für die Verwaltung und Verarbeitung der Daten wurden folgende Werkzeuge genutzt:
\begin{itemize}
    \item \textbf{GitHub:} Versionskontrolle und Zusammenarbeit im Team. Die Nutzung von Branches (z. B. \texttt{trial}-Branch) ermöglichte eine parallele Bearbeitung von Skripten und Datensätzen.
    \item \textbf{Python (Pandas):} Zentrale Bibliothek für die Datenmanipulation. Pandas wurde gewählt, da es eine effiziente Harmonisierung unterschiedlicher Tabellenstrukturen und eine automatisierte Berechnung neuer Metriken ermöglicht.
\end{itemize}

\section*{Vorgehen bei der Datenzusammenführung (Data Merging)}
Ein entscheidender Schritt war die Kombination der einzelnen Distrikt-Tabellen zu einem Gesamtdatensatz (\texttt{merged\_london\_data.csv}). Das Vorgehen gliederte sich in:
\begin{enumerate}
    \item \textbf{Standardisierung:} Vereinheitlichung der Spaltennamen, da die Originalscans teilweise leicht variierende Schreibweisen (z. B. "Heads of Famlies" vs. "Heads of Families") aufwiesen.
    \item \textbf{Konkatenation:} Zusammenführung der Daten aus den Boroughs Bethnal, Green, Hackney, Mile End Old Town, Poplar, Shoreditch, Stepney und St. George's-in-the-East.
    \item \textbf{Aggregatbildung:} Integration der Zusammenfassungen auf Borough-Ebene (\texttt{classes\_boroughs\_total.csv}) und Berufsebene (\texttt{professions\_total.csv}) zur Validierung der Gesamtsummen.
\end{enumerate}